\section{\texorpdfstring{Unbounded resolutions,\allowbreak{} derived adjunction,\allowbreak{} and products of K-injectives}{Unbounded  derived  and products of K-injectives}}\label{proofs-9206-9744-unbounded-resolutions-derived-adjunction-and-products-of-k-injectives}

Private source-order review of original derived.tex 9206–\allowbreak{}9744. Authority:\allowbreak{} Stacks commit a04446e5\allowbreak{}7ec1fbc2\allowbreak{}52a871af\allowbreak{}cec7752f\allowbreak{}b2807b14. Receiving public edition:\allowbreak{} 95a51593\allowbreak{}aceb247a\allowbreak{}c7a18311\allowbreak{}1678c3a0\allowbreak{}f8b8f4b5. The recorder binds the source files,\allowbreak{} original reports,\allowbreak{} operations,\allowbreak{} and exact proof identity. Editorial consequences remain separate from the source body. No novelty is claimed.

Human sources are \href{https://github.com/stacks/stacks-project/blob/a04446e57ec1fbc252a871afcec7752fb2807b14/derived.tex\#L9206}{the original Derived Categories TeX} and the referenced Categories and Homology chapters at the same authority. Secondary online comparisons located \href{https://stacks.math.columbia.edu/tag/0794}{the unbounded existence proposition} and \href{https://stacks.math.columbia.edu/tag/0BK6}{the K-injective product lemma}; the frozen TeX,\allowbreak{} rather than the mutable online edition,\allowbreak{} determines the corrections.

\subsection{1. A strict resolution system with its actual upper bounds}\label{proofs-9206-9744-a-strict-resolution-system-with-its-actual-upper-bounds}

Let A be abelian and let P contain zero,\allowbreak{} be closed under finite direct sums,\allowbreak{} and provide an epimorphism from an object of P onto each object of A. For an arbitrary complex K write T\_\allowbreak{}n\allowbreak{}=τ\_\allowbreak{}\{\textbackslash{}leq n\}K,\allowbreak{} n≥1. We construct termwise epimorphic quasi-isomorphisms p\_\allowbreak{}n:\allowbreak{}P\_\allowbreak{}n\allowbreak{}→T\_\allowbreak{}n with

P\_\allowbreak{}n\textasciicircum{}r\allowbreak{}=0 for r>n,\allowbreak{}

strictly compatible with termwise split monomorphisms P\_\allowbreak{}n\allowbreak{}→P\_\allowbreak{}\{n\allowbreak{}+1\},\allowbreak{} whose termwise cokernels belong to P.

The bounded resolution lemma,\allowbreak{} current derived.tex 5654–\allowbreak{}5737,\allowbreak{} gives P\_\allowbreak{}1\allowbreak{}→T\_\allowbreak{}1 with upper bound 1 and with all the required termwise properties. Its descending construction is worth specifying:\allowbreak{} to extend a partial resolution in degrees ≥r,\allowbreak{} choose an epimorphism from P\^{}\{r-1\} onto K\textasciicircum{}\{r-1\}×\_\allowbreak{}\{ker(d\_\allowbreak{}K\textasciicircum{}r)\allowbreak{}\}ker(d\_\allowbreak{}P\textasciicircum{}r)\allowbreak{}. Projection to the two factors gives the augmentation and differential. Surjectivity onto cycles at degree r and exactness on the previously established cohomology follow from this pullback. Starting with zero terms above the bound produces the stated resolution. This is the bounded construction used below.

Suppose P\_\allowbreak{}n and p\_\allowbreak{}n are given. Let Y\allowbreak{}=T\_\allowbreak{}\{n\allowbreak{}+1\} and let a:\allowbreak{}P\_\allowbreak{}n\allowbreak{}→Y be the composite with the canonical truncation inclusion. Use the actual source cone

C(a)\allowbreak{}\textasciicircum{}r\allowbreak{}=Y\textasciicircum{}r\allowbreak{}⊕P\_\allowbreak{}n\textasciicircum{}\{r\allowbreak{}+1\},\allowbreak{} d\_\allowbreak{}C\textasciicircum{}r\allowbreak{}=[[d\_\allowbreak{}Y\textasciicircum{}r,\allowbreak{}a\textasciicircum{}\{r\allowbreak{}+1\}]\allowbreak{},\allowbreak{}[0,\allowbreak{}-d\_\allowbreak{}\{P\_\allowbreak{}n\}\textasciicircum{}\{r\allowbreak{}+1\}]\allowbreak{}]\allowbreak{}.

It vanishes above n\allowbreak{}+1. Choose a termwise epimorphic quasi-isomorphism q:\allowbreak{}Q\allowbreak{}→C(a)\allowbreak{} from a complex with terms in P and Q\textasciicircum{}r\allowbreak{}=0 for r>n\allowbreak{}+1. Write q\textasciicircum{}r\allowbreak{}=(u\textasciicircum{}r,\allowbreak{}v\textasciicircum{}r)\allowbreak{},\allowbreak{} so

d\_\allowbreak{}Y u\allowbreak{}+a v\allowbreak{}=u d\_\allowbreak{}Q,\allowbreak{} -d\_\allowbreak{}\{P\_\allowbreak{}n\}v\allowbreak{}=v d\_\allowbreak{}Q. (1.1)\allowbreak{}

Define the next complex and map by

P\_\allowbreak{}\{n\allowbreak{}+1\}\textasciicircum{}r\allowbreak{}=P\_\allowbreak{}n\textasciicircum{}r\allowbreak{}⊕Q\textasciicircum{}r,\allowbreak{} d\_\allowbreak{}\{n\allowbreak{}+1\}\textasciicircum{}r\allowbreak{}=[[d\_\allowbreak{}\{P\_\allowbreak{}n\}\textasciicircum{}r,\allowbreak{}-v\textasciicircum{}r]\allowbreak{},\allowbreak{}[0,\allowbreak{}d\_\allowbreak{}Q\textasciicircum{}r]\allowbreak{}]\allowbreak{},\allowbreak{} p\_\allowbreak{}\{n\allowbreak{}+1\}\textasciicircum{}r\allowbreak{}=(a\textasciicircum{}r,\allowbreak{}u\textasciicircum{}r)\allowbreak{}:\allowbreak{}P\_\allowbreak{}n\textasciicircum{}r\allowbreak{}⊕Q\textasciicircum{}r\allowbreak{}→Y\textasciicircum{}r. (1.2)\allowbreak{}

Equation (1.1)\allowbreak{} gives d\_\allowbreak{}\{n\allowbreak{}+1\}²\allowbreak{}=0 and p\_\allowbreak{}\{n\allowbreak{}+1\}d\_\allowbreak{}\{n\allowbreak{}+1\}\allowbreak{}=d\_\allowbreak{}Y p\_\allowbreak{}\{n\allowbreak{}+1\}. The inclusion of P\_\allowbreak{}n is a chain map,\allowbreak{} and its composite with p\_\allowbreak{}\{n\allowbreak{}+1\} equals a strictly. The quotient complex is Q. Thus the inclusion is termwise split,\allowbreak{} its cokernel terms are in P,\allowbreak{} and the terms of P\_\allowbreak{}\{n\allowbreak{}+1\} belong to P. Its upper bound is n\allowbreak{}+1.

The map u\^{}r is the composite of the epimorphism q\^{}r with the biproduct projection C(a)\allowbreak{}\textasciicircum{}r\allowbreak{}→Y\textasciicircum{}r; hence u\textasciicircum{}r,\allowbreak{} and therefore p\_\allowbreak{}\{n\allowbreak{}+1\}\textasciicircum{}r,\allowbreak{} are epimorphisms. Finally,\allowbreak{} both C(p\_\allowbreak{}\{n\allowbreak{}+1\})\allowbreak{}\textasciicircum{}r and C(q)\allowbreak{}\textasciicircum{}r are,\allowbreak{} in the same ordered coordinates,\allowbreak{}

Y\textasciicircum{}r\allowbreak{}⊕P\_\allowbreak{}n\textasciicircum{}\{r\allowbreak{}+1\}\allowbreak{}⊕Q\textasciicircum{}\{r\allowbreak{}+1\},\allowbreak{}

and both differentials have the full matrix

[[d\_\allowbreak{}Y\textasciicircum{}r,\allowbreak{} a\textasciicircum{}\{r\allowbreak{}+1\},\allowbreak{} u\textasciicircum{}\{r\allowbreak{}+1\}]\allowbreak{},\allowbreak{} [0,\allowbreak{} -d\_\allowbreak{}\{P\_\allowbreak{}n\}\textasciicircum{}\{r\allowbreak{}+1\},\allowbreak{} v\textasciicircum{}\{r\allowbreak{}+1\}]\allowbreak{},\allowbreak{} [0,\allowbreak{} 0,\allowbreak{} -d\_\allowbreak{}Q\textasciicircum{}\{r\allowbreak{}+1\}]\allowbreak{}]\allowbreak{}.

The associativity map of these biproducts is an isomorphism of complexes. Since q is a quasi-isomorphism,\allowbreak{} its cone is acyclic; the same is true of C(p\_\allowbreak{}\{n\allowbreak{}+1\})\allowbreak{}. Thus p\_\allowbreak{}\{n\allowbreak{}+1\} is a quasi-isomorphism. This proves the induction. No compatibility merely in the homotopy category has been substituted for compatibility of the actual resolution maps.

DERIVED-UNDER-019 records the term bounds P\_\allowbreak{}n\textasciicircum{}r\allowbreak{}=0 for r\textgreater n and this explicit strict construction,\allowbreak{} as well as its dual in Section 4. The source states only boundedness above,\allowbreak{} although its hypotheses permit these sharper bounds. The proof is editorial; the existing source proof is also valid as checked next.

\subsection{\texorpdfstring{2. The original homotopy adjustment,\allowbreak{} disk correction,\allowbreak{} and colimit map}{2. The original homotopy  disk  and colimit map}}\label{proofs-9206-9744-the-original-homotopy-adjustment-disk-correction-and-colimit-map}

The source first constructs a middle complex P\textquotesingle{} in a triangle and a map f:\allowbreak{}P'\allowbreak{}→Y with f i-a\allowbreak{}=d h\allowbreak{}+h d on P\_\allowbreak{}n. Here i:\allowbreak{}P\_\allowbreak{}n\allowbreak{}→P' is termwise split. Choose graded retractions r\textasciicircum{}m:\allowbreak{}P'\textasciicircum{}m\allowbreak{}→P\_\allowbreak{}n\textasciicircum{}m and extend the homotopy by h̃\textasciicircum{}m\allowbreak{}=h\textasciicircum{}m r\^{}m. Then

f̃\allowbreak{}=f-d\_\allowbreak{}Y h̃-h̃ d\_\allowbreak{}\{P'\}

is a chain map homotopic to f,\allowbreak{} and f̃ i\allowbreak{}=a strictly. The equality holds because i is a chain map and r i\allowbreak{}=id; r need not be a chain map. This verifies the source\textquotesingle s omitted adjustment,\allowbreak{} including its sign.

After that adjustment,\allowbreak{} the source adds the contractible disk D with copies of an object P in degrees n,\allowbreak{}n\allowbreak{}+1,\allowbreak{} differential id\_\allowbreak{}P,\allowbreak{} and map to T\_\allowbreak{}\{n\allowbreak{}+1\} given by q:\allowbreak{}P\allowbreak{}→K\textasciicircum{}n and d\_\allowbreak{}K\textasciicircum{}n q:\allowbreak{}P\allowbreak{}→ker(d\_\allowbreak{}K\textasciicircum{}\{n\allowbreak{}+1\})\allowbreak{}. Here q is epic. The map is a chain map,\allowbreak{} the disk is contractible,\allowbreak{} and adding it preserves the quasi-isomorphism and the split inclusion of P\_\allowbreak{}n. Surjectivity in degree n is immediate. In degree n\allowbreak{}+1,\allowbreak{} the disk image is im(d\_\allowbreak{}K\textasciicircum{}n)\allowbreak{}. The map H\textasciicircum{}\{n\allowbreak{}+1\}(f̃)\allowbreak{} is epic; hence the image of f̃ in ker(d\_\allowbreak{}K\textasciicircum{}\{n\allowbreak{}+1\})\allowbreak{},\allowbreak{} together with im(d\_\allowbreak{}K\textasciicircum{}n)\allowbreak{},\allowbreak{} is all of ker(d\_\allowbreak{}K\textasciicircum{}\{n\allowbreak{}+1\})\allowbreak{}. This proves termwise surjectivity in the remaining degree. At degrees below n it already follows from the previous epimorphic augmentation. Thus the disk argument is correct; Section 1 gives a choice for which this final disk is unnecessary.

Assume A has exact sequential colimits. Put P\_\allowbreak{}∞\allowbreak{}=colim\_\allowbreak{}n P\_\allowbreak{}n termwise. The compatible p\_\allowbreak{}n give a chain map p:\allowbreak{}P\_\allowbreak{}∞\allowbreak{}→K because colim\_\allowbreak{}n T\_\allowbreak{}n\allowbreak{}=K. Indeed in each fixed degree r the term and adjacent differential of T\_\allowbreak{}n are those of K for all sufficiently large n. Exactness of colimits,\allowbreak{} applied to kernels,\allowbreak{} images and their quotient sequences,\allowbreak{} gives

H\textasciicircum{}r(P\_\allowbreak{}∞)\allowbreak{}\allowbreak{}=colim\_\allowbreak{}n H\textasciicircum{}r(P\_\allowbreak{}n)\allowbreak{}\allowbreak{}=H\textasciicircum{}r(K)\allowbreak{}.

For n≥max(1,\allowbreak{}r)\allowbreak{},\allowbreak{} the maps identify H\textasciicircum{}r(P\_\allowbreak{}n)\allowbreak{} with H\textasciicircum{}r(K)\allowbreak{} and are identities under these identifications. Hence p is a quasi-isomorphism. It is termwise epic too:\allowbreak{} choose n>r,\allowbreak{} so P\_\allowbreak{}n\textasciicircum{}r\allowbreak{}→K\textasciicircum{}r is already epic,\allowbreak{} and factor it through P\_\allowbreak{}∞\textasciicircum{}r.

The upper bound from Section 1 means that the map P\_\allowbreak{}n\allowbreak{}→P\_\allowbreak{}∞ factors strictly through τ\_\allowbreak{}\{\textbackslash{}leq n\}P\_\allowbreak{}∞. At degree n its image lies in the kernel of d\_\allowbreak{}\{P\_\allowbreak{}∞\}\textasciicircum{}n because P\_\allowbreak{}n\textasciicircum{}\{n\allowbreak{}+1\}\allowbreak{}=0. This factorization is a quasi-isomorphism,\allowbreak{} by the preceding cohomology calculation. This supplies exactly the stage maps used in the source\textquotesingle s comparison proof; it does not assume that a cohomological upper bound alone gives a factorization of arbitrary chain maps.

For each fixed degree r,\allowbreak{} construction (1.2)\allowbreak{} gives the compatible decompositions

P\_\allowbreak{}n\textasciicircum{}r\allowbreak{}=P\_\allowbreak{}1\textasciicircum{}r\allowbreak{}⊕Q\_\allowbreak{}1\textasciicircum{}r\allowbreak{}⊕…\allowbreak{}⊕Q\_\allowbreak{}\{n-1\}\textasciicircum{}r,\allowbreak{} P\_\allowbreak{}∞\textasciicircum{}r\allowbreak{}=P\_\allowbreak{}1\textasciicircum{}r\allowbreak{}⊕\allowbreak{}⊕\_\allowbreak{}\{n≥1\}Q\_\allowbreak{}n\textasciicircum{}r. (2.1)\allowbreak{}

The last isomorphism follows from the coproduct universal property:\allowbreak{} a compatible cocone on the finite partial sums is exactly one map out of each listed summand. This observation concerns the terms; the complex differential retains all the off-diagonal maps -v\_\allowbreak{}n from (1.2)\allowbreak{}.

\subsection{3. Unbounded left derivation and the weaker additive hypothesis}\label{proofs-9206-9744-unbounded-left-derivation-and-the-weaker-additive-hypothesis}

Here is the complete strengthened assertion,\allowbreak{} DERIVED-UNDER-020. Let A,\allowbreak{}B be abelian categories with exact sequential colimits. Let P satisfy Section 1\textquotesingle s hypotheses,\allowbreak{} and let F:\allowbreak{}A\allowbreak{}→B be additive. Suppose:\allowbreak{}

* F takes every bounded above acyclic complex with terms in P to an acyclic complex; * F commutes with countable coproducts.

Then the termwise functor has a left derived functor on all of D(A)\allowbreak{}. For the system of Section 1 it is computed by F(P\_\allowbreak{}∞)\allowbreak{}. The original proposition follows:\allowbreak{} its right exact functor is additive,\allowbreak{} and preservation of sequential colimits implies preservation of countable coproducts by expressing them as colimits of finite partial sums. Right exactness is not used in the argument below.

First consider bounded above complexes. Each has a resolution by a bounded above complex with terms in P. A quasi-isomorphism between two such resolutions has an acyclic cone whose terms are finite sums of P objects. F preserves the actual cone matrix because it is additive. The first assumption therefore makes its image a quasi-isomorphism. The left version of the source\textquotesingle s existence-and-computation criterion (current 5444–\allowbreak{}5470)\allowbreak{} produces LF\^{}- on D\textasciicircum{}-(A)\allowbreak{},\allowbreak{} computed by these resolutions. Canonical upper truncation supplies bounded representatives of cohomologically bounded-above objects; all comparisons are in D\textasciicircum{}-(A)\allowbreak{}.

Let S be the class of complexes that admit a presentation as a colimit of Section 1\textquotesingle s systems,\allowbreak{} with the stage maps to their own truncations as in Section 2. Every complex receives a quasi-isomorphism from an object of S. For P,\allowbreak{}Q in S choose witnesses P\_\allowbreak{}n,\allowbreak{}Q\_\allowbreak{}n and let α:\allowbreak{}P\allowbreak{}→Q be any quasi-isomorphism. It need not preserve these witnesses.

Write e\_\allowbreak{}\{P,\allowbreak{}n\}:\allowbreak{}F(P\_\allowbreak{}n)\allowbreak{}\allowbreak{}→LF\textasciicircum{}-(τ\_\allowbreak{}\{\textbackslash{}leq n\}P)\allowbreak{} and e\_\allowbreak{}\{Q,\allowbreak{}n\}:\allowbreak{}F(Q\_\allowbreak{}n)\allowbreak{}\allowbreak{}→LF\textasciicircum{}-(τ\_\allowbreak{}\{\textbackslash{}leq n\}Q)\allowbreak{} for the isomorphisms induced by their resolution maps. The maps

θ\_\allowbreak{}n\allowbreak{}=e\_\allowbreak{}\{Q,\allowbreak{}n\}\textasciicircum{}\{-1\} LF\textasciicircum{}-(τ\_\allowbreak{}\{\textbackslash{}leq n\}α)\allowbreak{}e\_\allowbreak{}\{P,\allowbreak{}n\} :\allowbreak{}F(P\_\allowbreak{}n)\allowbreak{}\allowbreak{}→F(Q\_\allowbreak{}n)\allowbreak{} (3.1)\allowbreak{}

are isomorphisms in D(B)\allowbreak{}. They are compatible with the actual maps P\_\allowbreak{}n\allowbreak{}→P\_\allowbreak{}\{n\allowbreak{}+1\},\allowbreak{} Q\_\allowbreak{}n\allowbreak{}→Q\_\allowbreak{}\{n\allowbreak{}+1\}:\allowbreak{} the truncation inclusions commute with τ\_\allowbreak{}\{\textbackslash{}leq n\}α,\allowbreak{} LF\^{}- is a functor,\allowbreak{} and the identifications e respect maps between the P resolutions. This proves stage compatibility,\allowbreak{} not just existence of unrelated isomorphisms.

It remains to identify their colimit with the map induced by F(α)\allowbreak{}. In K\textasciicircum{}-(A)\allowbreak{},\allowbreak{} choose a common roof R\allowbreak{}→P\_\allowbreak{}n and R\allowbreak{}→Q\_\allowbreak{}n over τ\_\allowbreak{}\{\textbackslash{}leq n\}α,\allowbreak{} commuting up to homotopy after the augmentation maps. One obtains such a roof by the calculus of quasi-isomorphisms,\allowbreak{} then by truncating its numerator above n and resolving it by a bounded above complex of P objects. Both R\allowbreak{}→P\_\allowbreak{}n and R\allowbreak{}→Q\_\allowbreak{}n are quasi-isomorphisms. Denote them a and b. The bounded computation says θ\_\allowbreak{}n\allowbreak{}=F(b)\allowbreak{}F(a)\allowbreak{}\textasciicircum{}\{-1\}. After composing to P and Q,\allowbreak{} the roof identity gives the homotopy α i\_\allowbreak{}\{P,\allowbreak{}n\}a∼i\_\allowbreak{}\{Q,\allowbreak{}n\}b. Additivity of F preserves this homotopy,\allowbreak{} so in D(B)\allowbreak{}

F(α)\allowbreak{} F(i\_\allowbreak{}\{P,\allowbreak{}n\})\allowbreak{}\allowbreak{}=F(i\_\allowbreak{}\{Q,\allowbreak{}n\})\allowbreak{}θ\_\allowbreak{}n. (3.2)\allowbreak{}

The definition (3.1)\allowbreak{} proves independence of all roof choices. It also proves that equality of alternative roofs,\allowbreak{} which can be checked after a common refinement,\allowbreak{} is carried to equality,\allowbreak{} not only to an unspecified isomorphism.

By (2.1)\allowbreak{} and preservation of countable coproducts,\allowbreak{} F(P)\allowbreak{}\allowbreak{}=colim\_\allowbreak{}n F(P\_\allowbreak{}n)\allowbreak{} and F(Q)\allowbreak{}\allowbreak{}=colim\_\allowbreak{}n F(Q\_\allowbreak{}n)\allowbreak{},\allowbreak{} canonically degree by degree. The comparison is a chain map because it commutes with the maps from each finite stage and those maps determine a map out of the colimit. Exact sequential colimits in B therefore give

H\textasciicircum{}r(F(P)\allowbreak{})\allowbreak{}\allowbreak{}=colim\_\allowbreak{}n H\textasciicircum{}r(F(P\_\allowbreak{}n)\allowbreak{})\allowbreak{},\allowbreak{} H\textasciicircum{}r(F(Q)\allowbreak{})\allowbreak{}\allowbreak{}=colim\_\allowbreak{}n H\textasciicircum{}r(F(Q\_\allowbreak{}n)\allowbreak{})\allowbreak{}.

Equations (3.1)\allowbreak{} and (3.2)\allowbreak{} identify H\textasciicircum{}r(F(α)\allowbreak{})\allowbreak{} with the colimit of the compatible isomorphisms H\textasciicircum{}r(θ\_\allowbreak{}n)\allowbreak{},\allowbreak{} so it is an isomorphism for every r. Thus F sends every quasi-isomorphism between members of S to a quasi-isomorphism. The same existence-and-computation criterion,\allowbreak{} now in the full homotopy category,\allowbreak{} constructs LF everywhere,\allowbreak{} with every member of S computing it. This proves the strengthened assertion.

We have retained exact sequential colimits in both categories. No claim that exact coproducts alone make all sequential colimits exact is needed. The weaker preservation condition on F suffices because the systems being used are split in each degree,\allowbreak{} not because arbitrary colimits were ignored.

For a concrete instance beyond right exact functors,\allowbreak{} take abelian groups,\allowbreak{} P the free groups,\allowbreak{} and F(M)\allowbreak{}\allowbreak{}=M[ℓ]\allowbreak{}\allowbreak{}=ker(ℓ:\allowbreak{}M\allowbreak{}→M)\allowbreak{},\allowbreak{} for a prime ℓ. F is additive and commutes with countable coproducts,\allowbreak{} since an element of a direct sum is killed by ℓ exactly when each of its finitely many nonzero components is. It is zero on free groups,\allowbreak{} so it sends all the required bounded above P complexes to the zero complex. It is not right exact:\allowbreak{} the epimorphism Z\allowbreak{}→Z/\allowbreak{}ℓ induces 0\allowbreak{}→Z/\allowbreak{}ℓ. Nevertheless the preceding proof constructs LF\allowbreak{}=0 on all of D(Ab)\allowbreak{},\allowbreak{} since every term in (2.1)\allowbreak{} is free. This is its left derived functor,\allowbreak{} not its right derived functor. Without right exactness there is no claim that H\textasciicircum{}0(LF(M[0]\allowbreak{})\allowbreak{})\allowbreak{} equals F(M)\allowbreak{}.

The received direction corrections are now semantically reconciled. MC-STK-ERR-0876 makes α run from P to Q,\allowbreak{} as (3.1)\allowbreak{} and (3.2)\allowbreak{} require,\allowbreak{} and writes its resulting map as H\textasciicircum{}r(F(α)\allowbreak{})\allowbreak{}. P02-E0070 proposes reversing the arrow rather than exchanging the displayed objects; both yield the same typed map. The three admitted operations are sufficient. P02-E0071 and the French DERIVED-065 report are completed with this entire argument,\allowbreak{} including compatibility in n.

\subsection{4. The dual system with the exact reversed indexing}\label{proofs-9206-9744-the-dual-system-with-the-exact-reversed-indexing}

Apply Section 1 in A\^{}op using the complex D(K)\allowbreak{}\textasciicircum{}r\allowbreak{}=(K\textasciicircum{}\{-r\})\allowbreak{}\textasciicircum{}op and d\_\allowbreak{}\{D(K)\allowbreak{}\}\textasciicircum{}r\allowbreak{}=(d\_\allowbreak{}K\textasciicircum{}\{-r-1\})\allowbreak{}\textasciicircum{}op. This is a complex because its two consecutive differentials compose to the opposite of d\_\allowbreak{}K². It sends a chain map K\allowbreak{}→L to D(L)\allowbreak{}\allowbreak{}→D(K)\allowbreak{},\allowbreak{} and H\textasciicircum{}r(D(K)\allowbreak{})\allowbreak{}\allowbreak{}=H\textasciicircum{}\{-r\}(K)\allowbreak{}\textasciicircum{}op:\allowbreak{} kernels and cokernels exchange under passage to the opposite category. Direct sums become finite products,\allowbreak{} and epimorphisms become monomorphisms. The actual canonical truncations satisfy

D(τ\_\allowbreak{}\{\textbackslash{}geq-n\}K)\allowbreak{}\allowbreak{}=τ\_\allowbreak{}\{\textbackslash{}leq n\}D(K)\allowbreak{}.

At the endpoint this equality identifies the kernel in A\^{}op with the cokernel of d\_\allowbreak{}K\textasciicircum{}\{-n-1\}; away from the endpoint it is the identity. Apply D back to the constructed diagram. It gives τ\_\allowbreak{}\{\textbackslash{}geq-n\}K\allowbreak{}→I\_\allowbreak{}n with termwise monomorphic quasi-isomorphisms,\allowbreak{} I\_\allowbreak{}n\textasciicircum{}r\allowbreak{}=0 for r<-n,\allowbreak{} and compatible termwise split epimorphisms I\_\allowbreak{}\{n\allowbreak{}+1\}\allowbreak{}→I\_\allowbreak{}n whose kernel terms belong to the given class I. All squares commute as chain maps. This proves the original inverse-system lemma and the lower-bound part of UNDER-019. No exact inverse-limit hypothesis has been asserted,\allowbreak{} and no conclusion about K\allowbreak{}→lim I\_\allowbreak{}n is deduced from this construction alone.

\subsection{5. Derived adjunction with explicit representatives and naturality}\label{proofs-9206-9744-derived-adjunction-with-explicit-representatives-and-naturality}

Retain the source\textquotesingle s exact functors F:\allowbreak{}D\allowbreak{}→D' and G:\allowbreak{}D'\allowbreak{}→D with G left adjoint to F,\allowbreak{} and the multiplicative systems S,\allowbreak{}S'. Fix K and M at which RF and LG respectively exist. Let s:\allowbreak{}K\allowbreak{}→I range over S,\allowbreak{} and let t:\allowbreak{}P\allowbreak{}→M range over S\textquotesingle. The values F(I)\allowbreak{} are an essentially constant ind-system in (S')\allowbreak{}\textasciicircum{}\{-1\}D' with value RF(K)\allowbreak{},\allowbreak{} and G(P)\allowbreak{} form an essentially constant pro-system in S\^{}\{-1\}D with value LG(M)\allowbreak{}.

Write c\_\allowbreak{}s:\allowbreak{}F(I)\allowbreak{}\allowbreak{}→RF(K)\allowbreak{} for the ind-system structure maps and d\_\allowbreak{}t:\allowbreak{}LG(M)\allowbreak{}\allowbreak{}→G(P)\allowbreak{} for the pro-system structure maps,\allowbreak{} in the respective localized categories. The Hom formulas for essentially constant systems and for the calculus of fractions identify each side of

Hom\_\allowbreak{}\{(S')\allowbreak{}\textasciicircum{}\{-1\}D'\}(M,\allowbreak{}RF(K)\allowbreak{})\allowbreak{} ≅ Hom\_\allowbreak{}\{S\textasciicircum{}\{-1\}D\}(LG(M)\allowbreak{},\allowbreak{}K)\allowbreak{} (5.1)\allowbreak{}

with the filtered colimit,\allowbreak{} in the two indices s and t,\allowbreak{} of Hom\_\allowbreak{}\{D'\}(P,\allowbreak{}F(I)\allowbreak{})\allowbreak{}≅Hom\_\allowbreak{}D(G(P)\allowbreak{},\allowbreak{}I)\allowbreak{}. The first index is K/\allowbreak{}S and the second is (S'/\allowbreak{}M)\allowbreak{}\textasciicircum{}op. Thus they are independent filtered indices; changing the order of their colimits is justified by the universal property of a cocone on their product,\allowbreak{} not by an exchange of a limit with a colimit.

More concretely,\allowbreak{} if h:\allowbreak{}P\allowbreak{}→F(I)\allowbreak{} and h\textasciicircum{}{\Rsixtyonesymbolfont ♭}:\allowbreak{}G(P)\allowbreak{}\allowbreak{}→I are adjoint,\allowbreak{} then (5.1)\allowbreak{} sends

c\_\allowbreak{}s h t\^{}\{-1\} to s\textasciicircum{}\{-1\}h\textasciicircum{}{\Rsixtyonesymbolfont ♭}d\_\allowbreak{}t. (5.2)\allowbreak{}

Changing s or t along a refinement gives the same localized maps:\allowbreak{} the naturality identities of the original adjunction identify the refined h\textasciicircum{}{\Rsixtyonesymbolfont ♭},\allowbreak{} and the structure maps c\_\allowbreak{}s,\allowbreak{}d\_\allowbreak{}t are compatible with those refinements. Conversely any arrow on the right is represented,\allowbreak{} by essential constancy of the pro-system,\allowbreak{} by some map G(P)\allowbreak{}\allowbreak{}→K in the localization. A left fraction writes it as s\^{}\{-1\}b with b:\allowbreak{}G(P)\allowbreak{}\allowbreak{}→I. Adjoint b back to h:\allowbreak{}P\allowbreak{}→F(I)\allowbreak{} and use the left expression in (5.2)\allowbreak{}. Equality of fractions or of representatives is exactly equality after a refinement,\allowbreak{} so these constructions are inverse.

These formulas prove naturality for actual arrows in K and M. A localized arrow is a composite of such arrows and inverses of arrows in S or S\textquotesingle. Naturality for those invertible arrows implies naturality for their inverses by multiplying the naturality square by the inverse isomorphisms. Hence (5.1)\allowbreak{} is functorial on the full domains where the two derived functors are defined,\allowbreak{} as stated by the source.

MC-STK-ERR-0877 correctly replaces the two undefined abelian derived categories in the general proof by (S')\allowbreak{}\textasciicircum{}\{-1\}D' and S\^{}\{-1\}D. The two P02 reports and French DERIVED-066 are reconciled by (5.1)\allowbreak{}–\allowbreak{}(5.2)\allowbreak{}.

For adjoint functors between abelian categories,\allowbreak{} G preserves zero objects and finite coproducts,\allowbreak{} while F preserves zero objects and finite products. Since the latter are biproducts,\allowbreak{} both functors are additive. Their degreewise adjunction restricts to chain maps:\allowbreak{} naturality in each variable equates the two chain-map equations. It also equates the homotopy relations,\allowbreak{} by additivity and the same naturality for their degree -1 components. Thus it descends to an adjunction on homotopy categories. Both induced functors preserve the actual shift and cone formulas,\allowbreak{} so they are exact there. This proves the abelian-category specialization.

If RF and LG are everywhere defined,\allowbreak{} define the unit η\_\allowbreak{}M:\allowbreak{}M\allowbreak{}→RF LG(M)\allowbreak{} as the preimage of id\_\allowbreak{}\{LG(M)\allowbreak{}\} under (5.1)\allowbreak{},\allowbreak{} and the counit ε\_\allowbreak{}K:\allowbreak{}LG RF(K)\allowbreak{}\allowbreak{}→K as the image of id\_\allowbreak{}\{RF(K)\allowbreak{}\}. Naturality in M gives ε\_\allowbreak{}\{LG(M)\allowbreak{}\} LG(η\_\allowbreak{}M)\allowbreak{}\allowbreak{}=id\_\allowbreak{}\{LG(M)\allowbreak{}\}. For the other identity,\allowbreak{} naturality in K sends the image of η\_\allowbreak{}\{RF(K)\allowbreak{}\} under (5.1)\allowbreak{} to ε\_\allowbreak{}K. But the image of id\_\allowbreak{}\{RF(K)\allowbreak{}\} is also ε\_\allowbreak{}K,\allowbreak{} and (5.1)\allowbreak{} is injective; therefore RF(ε\_\allowbreak{}K)\allowbreak{}η\_\allowbreak{}\{RF(K)\allowbreak{}\}\allowbreak{}=id\_\allowbreak{}\{RF(K)\allowbreak{}\}. This proves the actual adjunction,\allowbreak{} including its two triangle identities. No such composition is asserted in a partial domain unless its terms exist.

\subsection{\texorpdfstring{6. K-injectivity,\allowbreak{} localization,\allowbreak{} and bounded below complexes}{6.   and bounded below complexes}}\label{proofs-9206-9744-k-injectivity-localization-and-bounded-below-complexes}

The defining vanishing for maps from acyclic complexes also holds for all their shifts. If I is K-injective and u:\allowbreak{}M\allowbreak{}→N is a quasi-isomorphism,\allowbreak{} apply Hom\_\allowbreak{}K(-,\allowbreak{}I)\allowbreak{} to the cone triangle. The neighboring Hom groups of C(u)\allowbreak{} and its shifts vanish,\allowbreak{} so precomposition by u is bijective. Conversely,\allowbreak{} this bijectivity identifies a fraction f s\textasciicircum{}\{-1\}:\allowbreak{}N\allowbreak{}→I with the unique g:\allowbreak{}N\allowbreak{}→I in K satisfying g s\allowbreak{}=f. It is independent of the fraction because a common refinement identifies both such g. It is also injective:\allowbreak{} if g maps to zero in D,\allowbreak{} some quasi-isomorphism s makes g s zero in K,\allowbreak{} whence g\allowbreak{}=0 by that same bijectivity. Thus Hom\_\allowbreak{}K(N,\allowbreak{}I)\allowbreak{}\allowbreak{}→Hom\_\allowbreak{}D(N,\allowbreak{}I)\allowbreak{} is an isomorphism. If this last condition holds and N is acyclic,\allowbreak{} N is zero in D,\allowbreak{} so its Hom to I is zero in K. This proves all three equivalences.

For a distinguished triangle in K,\allowbreak{} apply Hom\_\allowbreak{}K(N,\allowbreak{}-)\allowbreak{} for acyclic N and each of its shifts. Its long exact sequence implies that if two terms are K-injective,\allowbreak{} the third is as well. Shifts of a K-injective complex are K-injective because Hom\_\allowbreak{}K(N,\allowbreak{}I[r]\allowbreak{})\allowbreak{}\allowbreak{}=Hom\_\allowbreak{}K(N[-r]\allowbreak{},\allowbreak{}I)\allowbreak{}.

For completeness let I have injective terms and vanish below a. Given f:\allowbreak{}N\allowbreak{}→I from an acyclic N,\allowbreak{} construct h\textasciicircum{}r:\allowbreak{}N\textasciicircum{}r\allowbreak{}→I\textasciicircum{}\{r-1\} with f\allowbreak{}=d\_\allowbreak{}I h\allowbreak{}+h d\_\allowbreak{}N. Set h\textasciicircum{}r\allowbreak{}=0 for r≤a. At r\allowbreak{}=a,\allowbreak{} f\^{}a kills im(d\_\allowbreak{}N\textasciicircum{}\{a-1\})\allowbreak{}\allowbreak{}=ker(d\_\allowbreak{}N\textasciicircum{}a)\allowbreak{},\allowbreak{} so it factors through im(d\_\allowbreak{}N\textasciicircum{}a)\allowbreak{}⊆N\textasciicircum{}\{a\allowbreak{}+1\}; injectivity of I\^{}a extends this factor to h\textasciicircum{}\{a\allowbreak{}+1\}. Inductively f\textasciicircum{}r-d\_\allowbreak{}I\textasciicircum{}\{r-1\}h\textasciicircum{}r kills ker(d\_\allowbreak{}N\textasciicircum{}r)\allowbreak{}:\allowbreak{} on the image of d\_\allowbreak{}N\textasciicircum{}\{r-1\} this follows from the chain equation for f and the preceding homotopy equation,\allowbreak{} with d\_\allowbreak{}I²\allowbreak{}=0. Extend again to obtain h\textasciicircum{}\{r\allowbreak{}+1\}. The construction proves every component of the homotopy equation. Hence f is zero in K and I is K-injective.

\subsection{7. The Hom-complex index correction and the product theorem}\label{proofs-9206-9744-the-hom-complex-index-correction-and-the-product-theorem}

For arbitrary complexes K,\allowbreak{}I define the full Hom complex of abelian groups by

C\textasciicircum{}r\allowbreak{}=∏\_\allowbreak{}\{b\allowbreak{}∈Z\}Hom\_\allowbreak{}A(K\textasciicircum{}b,\allowbreak{}I\textasciicircum{}\{b\allowbreak{}+r\})\allowbreak{},\allowbreak{} (d\_\allowbreak{}C\textasciicircum{}r f)\allowbreak{}\textasciicircum{}b\allowbreak{}=d\_\allowbreak{}I\textasciicircum{}\{b\allowbreak{}+r\}f\textasciicircum{}b-(-1)\allowbreak{}\textasciicircum{}r f\textasciicircum{}\{b\allowbreak{}+1\}d\_\allowbreak{}K\textasciicircum{}b. (7.1)\allowbreak{}

Expanding d\_\allowbreak{}C\textasciicircum{}\{r\allowbreak{}+1\}d\_\allowbreak{}C\textasciicircum{}r gives the two terms containing d\_\allowbreak{}I² and d\_\allowbreak{}K²,\allowbreak{} which vanish,\allowbreak{} and two terms containing d\_\allowbreak{}I f d\_\allowbreak{}K whose coefficients -(-1)\allowbreak{}\textasciicircum{}r and -(-1)\allowbreak{}\textasciicircum{}\{r\allowbreak{}+1\} cancel. Thus it is a complex. A degree r cocycle is exactly a chain map K\allowbreak{}→I[r]\allowbreak{},\allowbreak{} since the target differential is (-1)\allowbreak{}\textasciicircum{}r d\_\allowbreak{}I. If f\allowbreak{}=d\_\allowbreak{}C\textasciicircum{}\{r-1\}h,\allowbreak{} the corresponding chain homotopy has components (-1)\allowbreak{}\textasciicircum{}r h\textasciicircum{}b:\allowbreak{}K\textasciicircum{}b\allowbreak{}→I[r]\allowbreak{}\textasciicircum{}\{b-1\}. Indeed its boundary is

(-1)\allowbreak{}\textasciicircum{}r d\_\allowbreak{}I((-1)\allowbreak{}\textasciicircum{}r h)\allowbreak{}\allowbreak{}+((-1)\allowbreak{}\textasciicircum{}r h)\allowbreak{}d\_\allowbreak{}K \allowbreak{}= d\_\allowbreak{}I h\allowbreak{}+(-1)\allowbreak{}\textasciicircum{}r h d\_\allowbreak{}K\allowbreak{}=d\_\allowbreak{}C\textasciicircum{}\{r-1\}h.

Conversely any such homotopy arises this way. Therefore H\textasciicircum{}r(C)\allowbreak{}\allowbreak{}=Hom\_\allowbreak{}K(K,\allowbreak{}I[r]\allowbreak{})\allowbreak{},\allowbreak{} including H\textasciicircum{}0(C)\allowbreak{}\allowbreak{}=Hom\_\allowbreak{}K(K,\allowbreak{}I)\allowbreak{}.

The six factors in original 9712–\allowbreak{}9714 and 9720–\allowbreak{}9722 have K\^{}\{-b\} where (7.1)\allowbreak{} requires K\^{}b. DERIVED-NEW-032 corrects them and makes the full Hom complex explicit. This is not an innocuous dummy-index choice:\allowbreak{} changing b to -b in all factors would also change the exponents of I. For example,\allowbreak{} take K\allowbreak{}=I\allowbreak{}=Q[-1]\allowbreak{} in complexes of Q-vector spaces,\allowbreak{} so both are Q in degree 1. Their identity gives Hom\_\allowbreak{}K(K,\allowbreak{}I)\allowbreak{}\allowbreak{}=Q. In the printed middle product,\allowbreak{} nonzero K\^{}\{-b\} requires b\allowbreak{}=-1,\allowbreak{} whereas nonzero I\^{}b requires b\allowbreak{}=1. The product is zero. In this example the two adjacent printed groups are zero too,\allowbreak{} so no differential could repair the asserted middle cohomology. I is a bounded below complex of injectives,\allowbreak{} so this example is within the stated lemma\textquotesingle s setting.

Now let \{I\_\allowbreak{}t\} be the given set-indexed family,\allowbreak{} with the termwise products I\textasciicircum{}m\allowbreak{}=∏\_\allowbreak{}t I\_\allowbreak{}t\textasciicircum{}m existing in A. The product universal property defines d\_\allowbreak{}I by its components d\_\allowbreak{}\{I\_\allowbreak{}t\}; its square is zero because all its projections are zero. Formula (7.1)\allowbreak{} gives the actual isomorphism of complexes

Hom\textasciicircum{}•(K,\allowbreak{}I)\allowbreak{} ≅ ∏\_\allowbreak{}t Hom\textasciicircum{}•(K,\allowbreak{}I\_\allowbreak{}t)\allowbreak{}.

In degree r this is the canonical rearrangement ∏\_\allowbreak{}b Hom(K\textasciicircum{}b,\allowbreak{}∏\_\allowbreak{}t I\_\allowbreak{}t\textasciicircum{}\{b\allowbreak{}+r\})\allowbreak{}≅∏\_\allowbreak{}t∏\_\allowbreak{}b Hom(K\textasciicircum{}b,\allowbreak{}I\_\allowbreak{}t\textasciicircum{}\{b\allowbreak{}+r\})\allowbreak{}. Both terms in the differential agree after each projection,\allowbreak{} so this is an isomorphism of complexes,\allowbreak{} not just of their graded objects. Products of abelian groups are exact:\allowbreak{} kernels are coordinatewise,\allowbreak{} and a family of surjections is surjective by choosing a preimage in each coordinate. Hence cohomology commutes with this product. In particular,\allowbreak{} for every K and r,\allowbreak{}

Hom\_\allowbreak{}K(K,\allowbreak{}I[r]\allowbreak{})\allowbreak{} ≅ ∏\_\allowbreak{}t Hom\_\allowbreak{}K(K,\allowbreak{}I\_\allowbreak{}t[r]\allowbreak{})\allowbreak{}. (7.2)\allowbreak{}

For acyclic K each right-hand group vanishes because I\_\allowbreak{}t is K-injective and K[-r]\allowbreak{} is acyclic. Thus the full Hom complexes are acyclic,\allowbreak{} which justifies the source\textquotesingle s acyclicity step at every degree,\allowbreak{} not only at zero. Equation (7.2)\allowbreak{} proves I is K-injective. Applying Section 6 to I and the I\_\allowbreak{}t,\allowbreak{} now for arbitrary K,\allowbreak{} yields

Hom\_\allowbreak{}D(K,\allowbreak{}I)\allowbreak{} ≅ ∏\_\allowbreak{}t Hom\_\allowbreak{}D(K,\allowbreak{}I\_\allowbreak{}t)\allowbreak{}.

These isomorphisms are induced by the actual product projections and are natural in K,\allowbreak{} so they prove the universal property of the product in D(A)\allowbreak{}. Only the displayed term products in A were required. No exactness of products in A was assumed:\allowbreak{} the exact products used above are products of abelian groups. The theorem remains true.

The received article correction at 9709,\allowbreak{} MC-STK-ERR-0883,\allowbreak{} is already present. P02-E0074 and French DERIVED-072 are reconciled separately from the new Hom-index finding; neither report previously states it.

\subsection{8. Receiving statements and the review boundary}\label{proofs-9206-9744-receiving-statements-and-the-review-boundary}

The special direct-system uses in current more-algebra.tex 14989–\allowbreak{}15023 and cohomology.tex 6465–\allowbreak{}6495 require precisely the strict compatible epimorphic quasi-isomorphisms and split term inclusions proved in Sections 1–\allowbreak{}2. The improved bounds and the same colimit map apply to those choices of P. Their subsequent K-flatness arguments are separate uses of their cited lemmas; this bounded check does not replace those proofs.

The inverse system used in current cohomology.tex 10284–\allowbreak{}10310 has the strict maps and lower bounds of Section 4. Its later inverse-limit claim is not certified by the direct-system colimit calculation. The missing word before ``produce a diagram'' is routed as a reading observation for that chapter\textquotesingle s own report join,\allowbreak{} not counted as an additional finding here.

Current sites-cohomology.tex 9337–\allowbreak{}9350 invokes the unbounded criterion for a left adjoint g\_\allowbreak{}!,\allowbreak{} which is right exact and preserves all colimits. The new additive criterion therefore preserves those verified hypotheses; the later verification of its acyclicity condition is outside the excerpt. The further cited reductions at 10439–\allowbreak{}10450 and 10836–\allowbreak{}10850 motivate the explicit compatible-stage proof of Section 3. These excerpts do not certify the entirety of either receiving theorem.

Current injectives.tex 2038–\allowbreak{}2042 constructs a derived product from K-injective resolutions. Section 7 supplies the corrected Hom calculation for that exact invocation. The product projections and derived universal property are unchanged. No new receiving source operation is needed.

Eight more physical reports are complete in groups DERIVED-RECON-070,\allowbreak{} 071 and 072. All eight correspond to the three existing units MC-STK-ERR-0876,\allowbreak{} 0877 and 0883,\allowbreak{} comprising six already applied operations. Totals:\allowbreak{} 125/\allowbreak{}176 reports,\allowbreak{} 72 groups,\allowbreak{} 116 already corrected reports,\allowbreak{} eight missing copyedit/\allowbreak{}notation reports and one missing mathematical-index report; 67 prior Derived units and 84 prior operations reconciled. NEW-032 adds four proposed operations,\allowbreak{} giving 61 cumulative operations. UNDER-019 and UNDER-020 bring the editorial consequence count to twenty. Semantic review ends at 9744. Sequential reading has reached 9760; the next incomplete lemma starts at 9746 and unread source begins at 9761. The overall goal and received-work priority remain active.
